\begin{pastexam}[label={px:ntut-114-ee-ds-01}]{北科大電機所（戊組） 114 學年度 第 I-1 題}
（單選；4 分）

\textbf{I.\quad Single-select Multiple Choice:}

Which of the following NOT typically implemented using a stack data structure? \textbf{(4\%)}
\begin{enumerate}[label=\alph*)]
\item Undo/Redo operations in a text editor.
\item Queue management in a restaurant waiting system.
\item Function call management in program execution.
\item Syntax checking for matching parentheses in a compiler.
\end{enumerate}
\end{pastexam}

\begin{pastsoln}
答案是 b)。

\begin{itemize}
\item a) 用堆疊：每做一個動作就 push，Undo 時 pop 出最近的那一個（後做的先復原）；Redo 再用另一個
堆疊存被復原的動作。
\item b) \textbf{不用堆疊}：餐廳排隊是先來先服務，先到的人先帶位，是佇列（FIFO，資結書第~6.2~節），
不是後進先出。
\item c) 用堆疊：函式呼叫時把返回位址、參數、區域變數 push 成一個堆疊框，返回時 pop，最後呼叫的
最先返回（資結書第~6.4~節）。
\item d) 用堆疊：左括號 push，遇到右括號就 pop 出最近的左括號比對（資結書第~6.1.4~節）。
\end{itemize}
\end{pastsoln}
